\section{Introduction}
\label{sec:introduction}

\ifdefined\preprintfrontierfigure
\preprintfrontierfigure
\fi

Block diffusion language models, or BDLMs, offer a promising approach to
combining strong language-model quality with efficient inference. They retain
autoregressive dependencies across consecutive token blocks while generating
the tokens within each block through parallel denoising. This combines the
causal structure of autoregressive models with the parallelism of
diffusion-based generation
\citep{arriola2025blockdiffusion,odonoghue2026diffusiongemma,fu2026nemotronlabsdiffusion}.

Coding agents and other long-running assistants accumulate interaction
histories~\citep{yang2024sweagent,yao2023react}, motivating long-context BDLM
training.
For these workloads, activation memory and long-context attention are major
bottlenecks. Existing long-context training systems commonly use context
parallelism, which partitions sequence activations across devices so that each
device stores and processes only part of the context. During attention, devices
exchange information across these partitions, allowing local queries to
incorporate relevant information from the full
sequence~\citep{liu2023ringattention,jacobs2023ulysses}.

Conventional context parallelism (CP) is poorly matched to BDLM training
because it partitions the combined clean-plus-corrupted representation by
token position and performs distributed attention across those position
shards. Clean K/V from earlier blocks can condition multiple later blocks,
whereas corrupted K/V for block $b$ are used only by queries from that block.
At every layer, CP therefore exchanges both clean and corrupted K/V during
forward and their corresponding gradients during backward. The corrupted K/V
consume cross-rank bandwidth despite having no reuse across blocks
(Figure~\ref{fig:block-parallel-overview}(a)).

\begin{figure}[t]
    \centering
    \resizebox{\linewidth}{!}{\definecolor{BPBlue}{HTML}{356F9F}
\definecolor{BPOrange}{HTML}{B66718}

\begin{tikzpicture}[
    x=1cm,
    y=1cm,
    font=\sffamily\tiny,
    panel/.style={draw=black!24, rounded corners=2.2pt, line width=0.6pt,
                  minimum width=4.60cm, minimum height=4.35cm},
    clean/.style={draw=BPBlue, fill=BPBlue!13, rounded corners=1.2pt,
                  line width=0.65pt, minimum height=0.48cm,
                  inner xsep=0.05cm, inner ysep=0.03cm, align=center},
    corrupt/.style={draw=BPOrange, fill=BPOrange!15, rounded corners=1.2pt,
                    dash pattern=on 2.3pt off 1.2pt, line width=0.72pt,
                    minimum height=0.48cm, inner xsep=0.05cm,
                    inner ysep=0.03cm, align=center},
    ranklabel/.style={draw=black!42, fill=black!4, rounded corners=1.5pt,
                     line width=0.55pt, minimum width=0.86cm,
                     minimum height=0.48cm, inner xsep=0.05cm,
                     inner ysep=0.03cm, align=center,
                     font=\sffamily\tiny\bfseries},
    rankframe/.style={draw=black!48, fill=black!2, rounded corners=1.8pt,
                     line width=0.62pt, minimum width=1.28cm,
                     minimum height=2.20cm},
    summary/.style={draw=black!32, fill=black!3, rounded corners=1.7pt,
                    line width=0.55pt, minimum width=4.12cm,
                    minimum height=0.82cm, text width=3.88cm,
                    inner xsep=0.08cm, inner ysep=0.04cm, align=center}
]

% Equal panel frames keep the ownership progression visually comparable.
\node[panel] at (2.40,2.18) {};
\node[panel] at (7.20,2.18) {};
\node[panel] at (12.00,2.18) {};

\node[font=\sffamily\scriptsize\bfseries, text=black!82]
      at (2.40,4.08) {(a) CP: token-wise sharding};
\node[font=\sffamily\scriptsize\bfseries, text=black!82]
      at (7.20,4.08) {(b) BP: block ownership};
\node[font=\sffamily\scriptsize\bfseries, text=black!82, align=center]
      at (12.00,4.02) {(c) CSBP: sharded context\\+ block ownership};

% (a) Conventional CP transports both streams, regardless of cross-block reuse.
\foreach \x/\r in {0.90/1,2.40/2,3.90/3} {
    \node[rankframe] at (\x,2.18) {};
    \node[font=\sffamily\tiny\bfseries, text=black!76]
          at (\x,3.00) {Rank \r};
    \node[text=BPBlue!88!black] at (\x,2.76) {Clean K/V};
    \node[text=BPOrange!90!black, align=center]
          at (\x,1.99) {Corrupted\\K/V};
    % Individual cells denote token positions, not denoising blocks.
    \foreach \i in {1,2,3} {
        \pgfmathtruncatemacro{\tokenpos}{3*(\r-1)+\i}
        \node[draw=BPBlue, fill=BPBlue!13, line width=0.65pt,
              minimum width=0.30cm, minimum height=0.40cm, inner sep=0pt]
              at ({\x+0.34*(\i-2)},2.46) {\tokenpos};
        \node[draw=BPOrange, fill=BPOrange!15, line width=0.65pt,
              minimum width=0.30cm, minimum height=0.40cm, inner sep=0pt]
              at ({\x+0.34*(\i-2)},1.57) {\tokenpos};
    }
}
\foreach \x in {1.47,2.97} {
    \draw[<->, >=stealth, draw=BPBlue, line width=0.8pt]
          (\x,2.46) -- (\x+0.36,2.46);
    \draw[<->, >=stealth, draw=BPOrange, line width=0.8pt]
          (\x,1.57) -- (\x+0.36,1.57);
}
\node[font=\sffamily\tiny\bfseries, text=black!72]
      at (2.40,3.46) {Small cells denote token positions};
\node[summary, font=\sffamily\tiny\bfseries, text=black!72]
      at (2.40,0.50) {Clean and corrupted K/V\\cross ranks during attention};

% (b) Pure BP: target blocks are complete, but clean prefixes are replicated.
\node[ranklabel] at (5.76,2.94) {Rank 1};
\node[ranklabel] at (5.76,2.23) {Rank 2};
\node[ranklabel] at (5.76,1.52) {Rank 3};

\node[corrupt, minimum width=0.86cm, font=\sffamily\tiny\bfseries]
      at (6.72,2.94) {Block 1};

\node[clean, minimum width=0.86cm] at (6.72,2.23) {Clean 1};
\node[corrupt, minimum width=0.86cm, font=\sffamily\tiny\bfseries]
      at (7.68,2.23) {Block 2};

\node[clean, minimum width=0.86cm] at (6.72,1.52) {Clean 1};
\node[clean, minimum width=0.86cm] at (7.68,1.52) {Clean 2};
\node[corrupt, minimum width=0.86cm, font=\sffamily\tiny\bfseries]
      at (8.64,1.52) {Block 3};

\node[font=\sffamily\tiny\bfseries, text=BPOrange!90!black]
      at (7.20,3.46) {Each target block has one owner};
\node[font=\sffamily\tiny, text=BPBlue!88!black]
      at (7.20,1.14) {Overlapping clean prefixes are replicated};
\node[summary, draw=BPOrange!60, fill=BPOrange!5,
      font=\sffamily\tiny\bfseries, text=black!72]
      at (7.20,0.50) {Corrupted K/V stay local;\\clean prefixes are replicated};

% (c) CSBP: clean state is token-sharded; corrupted computation is block-owned.
\foreach \x/\r in {10.50/1,12.00/2,13.50/3} {
    \node[rankframe] at (\x,2.18) {};
    \node[font=\sffamily\tiny\bfseries, text=black!76]
          at (\x,3.00) {Rank \r};
    \node[text=BPBlue!88!black] at (\x,2.76) {Clean K/V};
    \foreach \i in {1,2,3} {
        \pgfmathtruncatemacro{\tokenpos}{3*(\r-1)+\i}
        \node[draw=BPBlue, fill=BPBlue!13, line width=0.65pt,
              minimum width=0.30cm, minimum height=0.40cm, inner sep=0pt]
              at ({\x+0.34*(\i-2)},2.46) {\tokenpos};
    }
    \node[corrupt, minimum width=1.10cm, minimum height=0.68cm,
          text width=0.94cm, inner xsep=0.04cm,
          font=\sffamily\tiny\bfseries] at (\x,1.57)
          {Complete\\block \r};
}

\foreach \x in {11.07,12.57} {
    \draw[<->, >=stealth, draw=BPBlue, line width=0.8pt]
          (\x,2.46) -- (\x+0.36,2.46);
}

\node[font=\sffamily\tiny\bfseries, text=BPBlue!88!black]
      at (12.00,3.46) {Shared clean sequence is token-sharded};
\node[summary, draw=BPOrange!60, fill=BPOrange!5,
      font=\sffamily\tiny\bfseries, text=black!72]
      at (12.00,0.50) {Only shared clean K/V\\cross ranks during attention};
\end{tikzpicture}
}
    \caption{(a) Conventional CP shards by token position and exchanges both clean and corrupted K/V across ranks,
    despite corrupted K/V having no reuse across blocks. (b) BP keeps
    corrupted K/V local to each block's owner, but replicates overlapping
    clean prefixes. (c) CSBP also shards the shared clean context, so only
    clean K/V and their gradients cross ranks during attention.
    Figure~\ref{fig:csbp} details the resulting tensor communication.}
    \label{fig:block-parallel-overview}
\end{figure}

Our first insight is that we can eliminate this block-specific communication by
exploiting the BDLM objective's separability across corrupted blocks. We
introduce block parallelism (BP), which makes target-block computation a new
distributed parallelism dimension. BP assigns each complete corrupted-block
computation to one rank. Different ranks execute different block computations
concurrently while each block's corrupted Q/K/V, activations, and gradients
remain local (Figure~\ref{fig:block-parallel-overview}(b)).

BP alone replicates increasingly long and overlapping clean prefixes across
ranks. To scale BP to long contexts, we introduce context-sharded
block parallelism (CSBP): CP shards the shared clean sequence across the same
ranks that own complete corrupted-block computations. Each clean K/V shard
serves all local blocks whose prefixes include it, and its gradients accumulate
on its owner. CSBP thus distributes clean computation and storage while keeping
corrupted K/V and their gradients local
(Figure~\ref{fig:block-parallel-overview}(c)).

We make four contributions.

\begin{itemize}
    \item \textbf{Block parallelism.} We introduce target-block computation as
    a new distributed parallelism dimension for BDLM training. BP uses
    per-block loss separability to assign each complete corrupted-block
    computation, including its forward pass, loss, and backward pass, to one
    rank. Different blocks execute concurrently while their corrupted K/V and
    gradients remain local, eliminating block-specific cross-rank communication.

    \item \textbf{Context-sharded block parallelism.} We introduce CSBP to
    scale block parallelism to long contexts without replicating the clean
    sequence. Over the same ranks, BP assigns complete corrupted-block
    computations while CP shards the shared clean-sequence computation. Each
    clean K/V shard serves every block whose prefix contains it, and its owner
    accumulates the backward contributions from all such blocks. CSBP removes
    corrupted K/V from cross-rank attention communication and avoids
    replicating overlapping clean prefixes.

    \item \textbf{Efficient long-context training.} Across six models,
    CSBP delivers $256$K-context speedups of
    \textbf{1.18--1.45$\boldsymbol{\times}$} for SFT and
    \textbf{1.27--1.33$\boldsymbol{\times}$} for AR-to-BDLM conversion over the best baselines,
    while matching or reducing peak HBM (Table~\ref{tab:main-256k-results}).
    The advantage grows with CP/BP degree and context length
    (Figures~\ref{fig:parallel-degree-scaling}--\ref{fig:sequence-length-scaling}):
    at 512K, it reaches \textbf{1.19$\boldsymbol{\times}$} on
    NemotronDiffusion 14B and \textbf{1.61$\boldsymbol{\times}$} on
    DiffusionGemma 26B-A4B.

    \item \textbf{Faster speculative decoding training and stronger agent performance.}
    CSBP accelerates DFlash2 drafter training by
    \textbf{2.48$\boldsymbol{\times}$} at 512K and
    \textbf{7.59$\boldsymbol{\times}$} at 1M. Under matched 12-hour
    DiffusionGemma SFT budgets, CSBP achieves higher pass rates at every trained
    checkpoint on SWE-bench Verified and Terminal-Bench Lite.
\end{itemize}
